\chapter*{How to read this book}
\addcontentsline{toc}{chapter}{How to read this book}
This book develops the mathematics behind \emph{Degree profiles of
uniform polynomial consequences}. It begins with actual algebraic proof
lines and ends with a self-contained account of the linear degree gap.
The intended reader is comfortable with polynomials, coefficient
vectors, linear combinations and spans. We recall kernels and quotient
spaces and introduce tensor products through explicit calculations.
Nullstellensatz proofs, Polynomial Calculus and the required label
actions are introduced here. No prior knowledge of modular
representation theory, homological algebra or of the earlier
complexity-gap paper is assumed. Understanding the conclusion is easier
than checking the complete proof: the common-witness construction in
Chapters~\ref{tx:ch:linear-layers}--\ref{tx:ch:prime} requires sustained
work with the displayed maps.

The central problem is
\[
                    \Gamma_n\vdash G^{(n)},
\]
where both the axioms and the polynomial expression to be derived come
from fixed rules as the universe size $n$ changes. A target with a
bounded-support index can be included as well. Deriving $1$ is the
important refutation special case, but it is not the only case.
Here $\Gamma_n$ is the set of input polynomials, each interpreted as
an equation equal to zero; $G^{(n)}$ is the target polynomial, whose equation
we want to derive. The symbol $\vdash$ means that an algebraic proof
exists. The two proof systems and their degree conventions are defined
in Chapter~\ref{tx:ch:consequences}.

\section*{The question which organises the whole proof}
Suppose a uniform system at size $n$ has a proof with degree at most
$d$. Can that degree budget also work at $n+q$, for a suitable
prime-power block $q$? For an invariant uniform target $G^{(n)}$, the
picture is
\[
 \Gamma_n\vdash_{\le d}G^{(n)}
 \quad\Longrightarrow\quad
 \Gamma_{n+q}\vdash_{\le d}G^{(n+q)}
 \quad\Longrightarrow\quad
 \Gamma_{n+2q}\vdash_{\le d}G^{(n+2q)}\quad\Longrightarrow\cdots.
\]
Here the proof system is fixed as either NS or ordinary PC. The
axioms and target follow their prescribed uniform rules at each size;
the larger proof may have a different list of lines. The resource
preserved by lifting is the degree bound for the whole proof. A
\emph{degree budget} means an upper bound on the degrees allowed in
the proof; it is not a bound on its length or on the universe size.
In this picture, \emph{invariant} means that every permutation of the
universe labels fixes the target polynomial. A target with distinguished
indices has a different invariance condition, defined in
Chapter~\ref{tx:ch:uniformity}.

This is the main question of the book. First we describe exactly
what has to be transported. Then we construct progressively stronger
lifts, measuring how large $n$ must be in relation to $d$. Finally
we reverse that numerical condition to obtain degree lower bounds.
General indexed targets use the same membership principle, with an
additional argument to detect their images; their individual vectors
need not be literal images under the invariant-target lift.

\section*{The linear-algebra picture behind lifting}
At a fixed degree budget $d$, ordinary PC is a calculation in a finite
vector space. Start from usable axioms, close under linear combinations,
and multiply a vector by a variable whenever that vector has degree at
most $d-1$. This gives a subspace $C_d(n)$. NS uses a different subspace
$I_d(n)$, spanned by the permitted monomial--axiom products. The target
question is membership in one of these spaces.

The difficult theorem relates the membership calculations at $n$ and
$n+q$, where $q=p^a$ is a suitably large power of the characteristic.
One exact operation deletes the extra block of labels. A related linear
operator lifts smaller consequences into the larger universe. The
construction has to preserve actual maps, intersections and images;
a count of dimensions would not suffice.

For a first reading of this transition, begin with the worked
multiplication calculation in Section~\ref{tx:sec:multiplication-registers}.
Then follow the four operations making one closure stage in
Section~\ref{tx:sec:grammar}. Section~\ref{tx:sec:target-reflection}
explains why preserving this calculation still leaves a separate
target question. Chapter~\ref{tx:chap:deletion} introduces the objects
in Theorem~\ref{tx:thm:deletion} one by one and gives a guide to each
clause. At that point the theorem is a precise statement of what
the later construction will prove, not a fact to accept without proof.

There are two different operations which should never be confused.
Multiplication by a variable takes degree-$d-1$ vectors to degree-$d$
vectors \emph{within one universe}. The lifting operator takes
polynomials from the smaller universe to the larger one \emph{without
increasing degree}. It is linear, but is not generally a homomorphism
of polynomial rings.

\section*{Why one proof gives a constant bound on a progression}
Fix one degree $d$ and one eligible starting size $n$. Lifting gives
proofs with that same budget at $n+q,n+2q,\ldots$. The absolute budget
stays $d$. Once that starting proof has been selected, $d$ is a
\emph{constant} for the whole progression, even if it was originally
described as, say, a small multiple of $\log n$. At the later universe sizes
we keep its numerical value; we do not recompute the logarithm.
Its ratio to universe size tends to zero:
\[
                         \frac{d}{n+\ell q}\longrightarrow0.
\]
The constant bound initially belongs to one arithmetic progression.
It is not yet a constant bound over all universe sizes. Two-way transfer
at every lower threshold shows exactly which finite minimum degree
recurs on that progression. The four types describe how those
progressions fit together.

The lower-bound argument reverses this observation. If a finite degree
cannot recur, it must have occurred before transfer became available.
The theorem makes that availability threshold linear in $d$. Solving
this inequality for $d$ gives the linear lower bound. In the fluctuating
case, a separate bound on the length of the period locates finite linear
peaks. These are different uses of the same transfer theorem.

\section*{The stages are improvements in how far we can lift}
The earlier refutation classification gives logarithmic NS and
double-logarithmic PC lower bounds~\cite{Riis,RiisLong}. The underlying
transfer estimates summarised in the LICS paper have exponential and
double-exponential starting universe sizes as functions of the fixed degree.
These are historical refutation estimates, obtained there from
substantial earlier work; they are not elementary lemmas reconstructed
in this book.

The constructions developed here make the starting universe size much smaller:
\begin{center}
\begin{tabular}{@{}>{\raggedright\arraybackslash}p{.43\linewidth}>{\raggedright\arraybackslash}p{.25\linewidth}>{\raggedright\arraybackslash}p{.22\linewidth}@{}}
\toprule
Lifting stage & Sufficient starting size & Degree lower bound\\
\midrule
Earlier NS refutation result & at most $2^{cd}$ & $\Omega(\log n)$\\
Earlier PC refutation result & at most $2^{2^{cd}}$ & $\Omega(\log\log n)$\\
Characteristic-two construction with fresh labels & $O(d^2)$ & $\Omega(\sqrt n)$\\
One retained pool, every positive characteristic & $O(d)$ & $\Omega(n)$\\
\bottomrule
\end{tabular}
\end{center}
These are degree lower bounds on the hard sizes specified by the
classification, not universal degree upper bounds.
The constants depend on the fixed family and field. A small degree
means small enough to meet the relevant starting-size inequality.
Chapter~\ref{tx:chap:lift} gives an elementary counting explanation
for the prime-power steps. Chapter~\ref{tx:ch:quadratic} supplies the
first fully worked lifting construction. The central
advance in Chapters~\ref{tx:ch:linear-layers}--\ref{tx:ch:prime} is to
replace its quadratic label cost by a linear one, and to make that
construction work also in odd characteristic. The T4 conclusion uses
an additional linear bound on the period and gives finite linear peaks.

This development connects proof complexity with modular representation
theory and algebraic combinatorics. Those communities already study
representations built from bounded numbers of labels and periodic
behaviour as the universe grows~\cite{CEF,Harman,HKPSYZ,BDM}.
Section~\ref{tx:sec:stability-context} identifies the close antecedents.
The task here is to prove transport for the actual polynomial proof
rules and targets, with all the required finite constructions supplied
in the book.

\section*{A layered route through the construction}
The first part replaces proof syntax by consequence spaces, then works
through a Boolean first-order example over a fixed finite field: a
strict order in which every element chooses a smaller predecessor.
Chapter~\ref{tx:ch:running-order} proves the elementary growing-degree
obstruction and shows exactly how comparison and lifting strengthen it
to a linear bound. The predecessor may be many-to-one; the example
uses no matching or injection constraint. The second
part states the precise transfer contract and derives the four profiles
from it; it also identifies the desired lifting map. The third part
constructs the common algebraic projector which fulfils that contract,
first in characteristic two and then in every positive characteristic.
The fourth part puts the target and consequence maps into the same
argument and proves the independent NS--PC comparison. The final part
works through examples, dual witnesses, matching obstructions and the
logical undecidability consequence.

Within a technical chapter the order is always the same: say which
property is needed, explain why it is enough, prove the structural
lemma, and then inspect the calculation that realises it. The full
normalisation, zero-marginal, rook and odd-characteristic arguments
remain in the main chapters. Exercises have worked solutions at the
end, but the proofs of the main results do not depend on doing them.

A first reading can follow the contracts, examples and checkpoints.
A second reading can fill in each finite algebraic construction in turn.
This is a way of learning the complete proof, not permission to regard
its hard inputs as unproved assumptions.

When returning to a technical section, keep track of the space on
which an operator acts and which parameters are fixed. Powers of a
linear operator mean repeated composition. A polynomial degree, a
tuple length and a support budget measure different things, even
when the argument relates them by an inequality. Local reminders at
the main changes of construction, together with the notation guide
in the final chapter, identify these roles.

The most direct first route through the opening is to read the basic
template definition in Chapter~\ref{tx:ch:uniformity}, then the target
request and support budget in
Sections~\ref{tx:sec:target-request}--\ref{tx:sec:target-budgets}.
The orbit-sum and compression discussion can be read later. In
Chapter~\ref{tx:ch:closure}, follow the closure recurrence and its maps;
the coefficient-list size interpretation is a separate application.
Then follow the running order example in
Chapter~\ref{tx:ch:running-order}. Its final steps specify the two
general inputs still to be constructed: uniform comparison and
degree-preserving lifting.
Chapter~\ref{tx:chap:contract} then names the three obligations which
organise the rest of the proof: identify the data, preserve the
calculation, and reflect target vanishing. The chapters below refer
back to these obligations as each construction is supplied.

\section*{Specialist tools and a route through the proof}
The proof uses several specialist languages. Its representation-theoretic
part concerns permutation modules, group algebras, annihilator ideals
and nilpotent chains. Its transport language consists of tensor products,
quotients, naturality and exactness. Its finite constructions use common
partial inverses, ordered projectors, marginals, sections and rook
coordinates. Dual functionals and matching designs supply obstructions;
definability and effective reductions explain the logical consequences.
The background tools are developed at their first use and recalled
when their role changes. The following map identifies the main
transitions between mathematical subjects. It also helps a reader
returning to the proof locate the relevant definitions.

\begin{longtable}{@{}p{.23\linewidth}p{.47\linewidth}p{.22\linewidth}@{}}
\toprule
Tool & Concrete use in this book & First development\\
\midrule
Ideals and degree cutoffs & Polynomial combinations, with the
allowed products constrained by degree & Chapters~\ref{tx:ch:consequences}--\ref{tx:ch:closure}\\
Orbits and stabilisers & Grouping coefficient patterns under relabelling
& Chapter~\ref{tx:ch:uniformity}\\
Scalar extension & Solving the same coordinate problem with auxiliary
field coefficients & Chapter~\ref{tx:ch:closure}\\
Tensor products & Coordinates for paired inputs, with a diagonal
label action & Chapter~\ref{tx:chap:contract}\\
Naturality and exactness & Commuting maps and preserving the actual
kernel, image and quotient calculations & Chapter~\ref{tx:chap:contract}\\
Nilpotent chains & Reading which directions a block operator kills
or leaves in its quotient & Chapter~\ref{tx:chap:deletion}\\
Frobenius and sections & Prime-power counting and a chosen lift of
each retained vector & Chapter~\ref{tx:chap:lift}\\
Common partial inverses & One group-algebra expression certifying
the required projector on several spaces & Chapters~\ref{tx:chap:lift}--\ref{tx:ch:square-zero}\\
Marginals and layer sections & Compatible coordinate sums and
corrections using one retained label pool & Chapter~\ref{tx:ch:linear-layers}\\
Rook coordinates & Encoding overlap patterns and solving a
triangular coefficient system & Chapter~\ref{tx:ch:rooks}\\
Dual functionals and pairings & Coefficient tests separating targets
from consequence spaces & Chapters~\ref{tx:ch:running-order}--\ref{tx:chap:contract}\\
Matching designs & Bounded-degree coefficient tests and their
extension rules & Chapter~\ref{tx:app:matching}\\
Definability and reductions & The model-theoretic interpretation and
the transfer of an undecidable computation problem & Chapters~\ref{tx:sec:model-theory}--\ref{tx:sec:undecidability}\\
\bottomrule
\end{longtable}

\paragraph{Background reading.}
For ordinary algebra, Lang~\cite{LangAlgebra} has chapters on
modules, algebraic extensions, a single endomorphism, bilinear forms
and tensor products. Alperin~\cite{AlperinLocal} introduces the modular
representation theory of finite groups. The exactness and reflection
vocabulary is described in the Stacks Project~\cite{StacksExact}; here
we develop it through kernels and images of the displayed linear maps.
Cutland~\cite{Cutland}, especially the chapters on computable functions,
universal programs, undecidability and reducibility, gives background
for the counter-program reductions. The references in
Section~\ref{tx:sec:stability-context} locates the representation-stability
connection. Chapter~\ref{tx:sec:model-theory} explains the paper's
Section~11 topological classification, including a worked residue-tree
example. Section~\ref{tx:sec:euler-background} gives optional historical
background on Euler structures, not a criterion characterising T1--T4.
These are reading routes for the terminology. The common witness,
its natural map compatibility, the rook correction and the ordinary-PC
extension invariant are established explicitly in this book.

The opening route above is for learning the problem: it reaches the
running example and the transport contract quickly. The book's chapter
order is for learning the construction: it explains the destination
before constructing the engine, including the simpler quadratic stage.
The following dependency order is for checking that every input of the
linear theorem has been discharged; it returns to classification only
after transfer is proved. These purposes account for their different
orders. Each item
states a concrete conclusion to check, rather than asking the reader to
accept a general stability slogan.
\begin{enumerate}
\item \textbf{Replace proofs by linear algebra.}
In Chapters~\ref{tx:ch:consequences}--\ref{tx:ch:closure}, check that
$I_d$ and $C_d$ describe \emph{exactly} the legal degree-$d$ proofs.
In particular, PC intersects with $P_{\le d-1}$ after cancellation.
Check stable coefficients and the separate degree-zero convention.
\item \textbf{Identify the certificate for exactness.}
In Chapter~\ref{tx:ch:square-zero}, verify that one common identity
$zwz=z$ gives the projector $(1-zw)(1-wz)$ and exact deletion on
the entire common annihilator class. The two-point failure of exactness
explains why the moved-basis calculation alone is insufficient.
\item \textbf{Construct that common certificate in characteristic two.}
In Chapter~\ref{tx:ch:linear-layers}, check the partition of unity,
the signed layer sections and the correction interface
(Lemma~\ref{tx:lem:layer-interface}). In Chapter~\ref{tx:ch:rooks},
check strengthened forcing and triangular realisation. These prove that
one correction repairs length $b$ while annihilating every shorter
length, on the same pool of retained labels.
\item \textbf{Check the changes for an arbitrary prime.}
Chapter~\ref{tx:ch:prime} replaces $z$ by its top power
$u=z^{p-1}$ in the witness identity. Check purity on tuple generators,
the ordered projector factors, the $p$-level sections and the actual
transpose relabelling. Purity is not assumed for arbitrary nilpotent
representations.
\item \textbf{Return to the prescribed polynomials.}
In Chapter~\ref{tx:chap:assemble-transfer}, check the axiom and
multiplication maps entry by entry, then the full target map $J$.
For indexed targets, check the nonzero-subquotient pairing
(Lemma~\ref{tx:lem:detect}); exactness alone does not reflect vanishing.
The pairing doubles the original monomial support $kd$.
\item \textbf{Read off the lower bounds.}
Now check Chapter~\ref{tx:ch:profiles}. The linear onset controls
nonrecurring values in T2 and on T3 hard residues. The period bound
has the additional job of locating \emph{finite} T4 peaks. This step
is an argument about integer sequences after transfer has been proved.
\item \textbf{Check the separate comparison theorem.}
Chapter~\ref{tx:sec:comparison} uses a strict-chain bound and an NS
plateau to contain the whole PC closure. It proves same-type agreement;
it is not the source of the linear PC lower bound.
\end{enumerate}

Chapter~\ref{tx:ch:quadratic} is a complete alternative construction in
characteristic two, useful for seeing why spending fresh labels gives
only a square-root bound. The linear route does not depend on that
construction or its normalisation lemma. Chapters~\ref{tx:chap:deletion}
and~\ref{tx:chap:lift} state and explain the operation before its
coefficients are constructed; their promised input is discharged by
Chapters~\ref{tx:ch:square-zero}--\ref{tx:ch:prime}.
The later example, matching and undecidability chapters have their own
additional arguments.

The core formal results have independently checked Lean source builds,
theorem statements, axioms and dependencies, as recorded with the paper.
That evidence does not certify every new pedagogical formulation in this
book. A human check can follow the finite calculations above without
treating either the formalisation or an external representation-stability
theorem as a substitute for a missing step.

\section*{Notation which keeps the roles separate}
\begin{center}
\begin{tabular}{@{}ll@{}}
\toprule
Symbol & Meaning\\
\midrule
$p$ & Fixed prime characteristic of the coefficient field\\
$n$ & Retained, smaller universe size\\
$q=p^a$ & Size of the added/deleted block; $a$ is an exponent\\
$N=n+q$ & Larger universe size\\
$d$ & Degree bound for the whole proof\\
$G_b^{(n)},G^{(n)}$ & Indexed target and fully invariant target at size $n$\\
$D_G$ & Fixed upper bound on the target polynomial degree\\
$k,s,t$ & Supports of a variable, axiom index, target index\\
$r=R(d)$ & One support budget once $d$ has been fixed\\
$N(d)$ & Onset beyond which a degree-$d$ threshold transfers\\
$Q(d)$ & A suitable prime-power period for that threshold\\
$I_d(n),C_d(n)$ & NS and ordinary-PC consequence subspaces\\
$\mathsf L$ & Degree-preserving lift from $n$ to $N$\\
$\Del$ & Exact deletion operation from larger modules to smaller ones\\
\bottomrule
\end{tabular}
\end{center}
In particular, the support budget $r$ is \emph{not} the exponent $a$
in $p^a$, and the variable support $k$ is neither of them. The number
of polynomial variables can be of order $n^k$ and is not itself $n$.

\section*{Reading alongside the current paper}
This edition aligns terminology and target notation with the
78-page paper revised on 4 November 2026
(\href{https://github.com/SR123/degree-profiles-linear-bounds/blob/9e7cea1064f28f7e2ffd5eab44d52bfa854c5baa/paper/paper.pdf}{fixed GitHub version}). The paper
is the concise statement of the research; this book supplies the slower
route through its arguments. In particular $N(d)$ now denotes the
same periodicity onset in both texts. Ordinary $N=n+q$ still denotes
a larger universe size. The older refutation-only tutorials are
historical guides, rather than additional prerequisites.

\begin{center}\small
\begin{tabular}{@{}>{\raggedright\arraybackslash}p{.35\textwidth}>{\raggedright\arraybackslash}p{.59\textwidth}@{}}
\toprule
Question while reading the paper & Where to turn here\\
\midrule
What is uniform, and what does degree count? & Chapters 1--3; distinguish a full polynomial from its bounded support of distinguished labels.\\[3pt]
How does transfer imply the four cases? & Chapter~\ref{tx:ch:profiles}; compare \emph{thresholds}, then recurring exact values.\\[3pt]
Why does deletion preserve the proof? & Chapter~\ref{tx:chap:contract}; transport maps and intersections, not just dimensions.\\[3pt]
Where is the linear improvement? & Chapters~\ref{tx:ch:quadratic}--\ref{tx:ch:prime}; compare fresh labels with reuse of one common pool.\\[3pt]
Why is the prescribed target preserved? & Chapter~\ref{tx:chap:assemble-transfer}; the quotient obstruction must survive deletion.\\[3pt]
How does topology separate T1--T4? & Section~\ref{tx:sec:topological-spectrum}; the residue tree, compactness, and the binary worked example.\\[3pt]
What is formalised? & The accompanying Lean packet's scope and endpoint tables; the new topological reformulation has a written proof.\\
\bottomrule
\end{tabular}\end{center}

A short first pass is Chapters 1 and 3, the contract in
Chapter~\ref{tx:chap:contract}, Chapter~\ref{tx:ch:profiles}, and
Section~\ref{tx:sec:topological-spectrum}. Then return to the common
witness and rook constructions. Their purpose is to justify the
contract with a \emph{linear} label budget; the classification proof
itself is considerably simpler.

\tableofcontents
\mainmatter
